\cotefile{91}
\scanpage{129}
\mspage{- 135-}
% archive sheet number: 128
\hfill n\textsuperscript{o} 316

a) L'algèbre $\overset{\handwritten{\sim}}{U}$ et l'application $\handwritten{\unsure{\tilde{\varrho}|\tilde{\rho}}}$ sont solutions d'un problème d'application universelle : pour tout $p$-homomorphisme $f$ de $g$ dans une algèbre associative $B$ sur $K$ (considérée comme $p$-algèbre de Lie), il existe un $K$-homomorphisme et un seul $f'$ de $\overset{\handwritten{\sim}}{U}$ dans $B$ (pour les structures d'algèbres associatives) tel que $f=f'\unsure{o|\circ}\handwritten{\tilde{\varrho}}$.

b) Montrer que si $(b_{\handwritten{\lambda}})_{\handwritten{\lambda}\handwritten{\in}\handwritten{\Lambda}}$ est une base de $g$ ($\handwritten{\Lambda}$ étant totalement ordonné), $\handwritten{\tilde{\varrho}}$ est injective, et que, si on identifie $x$ et $\handwritten{\tilde{\varrho}}(x)$ pour $x\handwritten{\in} g$, les éléments $\handwritten{\prod_{\lambda}}b_{\handwritten{\lambda}}^{\handwritten{\nu_{\lambda}}}$ (où les $\handwritten{\lambda}$ vont en croissant, où les $\handwritten{\nu_{\lambda}}$ sont nuls sauf un nombre fini d'entre eux, et où $0\handwritten{\leqslant\nu_{\lambda}}\handwritten{<}p$ pour tout $\handwritten{\lambda}$) forment une base de $\overset{\handwritten{\sim}}{U}$. (Identifiant canoniquement $g$ à un sous-module de $U$, on pose $\phi(x)=x^{p}-x^{[p]}$ pour tout $x\handwritten{\in} g$. Pour tout indice composé $\alpha=(\alpha_{\handwritten{\lambda}})\handwritten{\in}\unsure{\mathbb{N}|N}^{(L)}$, soit $\alpha_{\handwritten{\lambda}}=\handwritten{\beta_{\lambda}}+p\handwritten{\gamma_{\lambda}}$ avec $0\handwritten{\leqslant\beta_{\lambda}}\handwritten{<}p$, et soit $T_{\alpha}=(\handwritten{\prod_{\lambda}}b_{\handwritten{\lambda}}^{\handwritten{\beta_{\lambda}}})(\handwritten{\prod_{\lambda}}(\phi(b_{\handwritten{\lambda}}))^{\handwritten{\gamma_{\lambda}}})$. Montrer que les $T_{\alpha}$ forment une base de $U$ et les $T_{\alpha}$ tels que $\handwritten{\gamma}=(\handwritten{\gamma_{\lambda}})\ne 0$ une base de $\unsure{J|\mathscr{J}}$ ; on observera que les $\phi(b_{\handwritten{\lambda}})$ appartiennent au centre de $U$).

7) - Soit $g$ une $p$-algèbre de Lie sur un anneau $K$ de caractéristique $p>0$ ($p$ premier). On dit qu'une dérivation $D$ de $g$ est une $p$\uline{-dérivation} si on a $D(x^{p})=(\operatorname{ad}x)^{p-1}.dx$ pour tout $x\handwritten{\in} g$. Toute dérivation intérieure est une $p$-dérivation.

a) Si $L$ est une $K$-algèbre associative, toute dérivation de $L$ est une $p$-dérivation lorsque $L$ est considérée comme $p$-algèbre de Lie. (Utiliser la formule (2) de l'exerc. 19 du \S 1).

b) Supposons que $g$ ait une base. Pour qu'une dérivation de $g$ soit une $p$-dérivation, il faut et il suffit qu'elle puisse se prolonger en une dérivation de l'algèbre enveloppante restreinte de $g$.
\scanpage{131}
% archive sheet number: 130
$V$ anneau à valuation discrète, corps résiduel $k$ parfait, corps des fractions $K$

Pour tout préschéma $X$ sur $V$, on définit $\bar X$ par Greenberg.
\[
\bar X(k')\simeq X(V_{k'})\deleted{\unsure{\longrightarrow|\longleftrightarrow}X}
\]
\margin{($k'$ alg. sur $k$, $V_{k'}$ alg. correspondante sur $V$).}

Si $k'$ est parfait, $V_{k'}$ est un anneau de valuation discrète, soit $K_{k'}=K'$ son corps des fractions, on a
\[
\bar X(\unsure{k|k'})\simeq X(V_{k'})\longrightarrow X(K')
\]

Soit \deleted{$\unsure{\varphi:}X\longrightarrow$} $Y$ un autre $V$ préschéma, \inserted{\uline{propre} \uline{sur} $V$,} et soit
\[
\varphi:\quad X_K\longrightarrow Y_K
\]
un morphisme de $K$-préschémas, nous allons en déduire une \uline{application ensembliste}
\[
\bar\varphi:\quad\bar X\longrightarrow\bar Y
\]

Soit en effet $\bar x\in\bar X$, soit $k'=\tilde{k}(\bar x)$ \deleted{une}\inserted{``la''} clôture parfaite de $k(\bar x)$, alors $\bar x$ définit un \unsure{pt|p\textsuperscript{t}} $\bar x'\in\bar X(k')$, d'où un pt \deleted{\unsure{dans} $\unsure{\varphi} + x$} $\alpha_X(\unsure{\bar x''|\bar x'})\in\bar X(V_{k'})\subset\bar X(K_{k'})$ et appliquant $\varphi$ un point
\[
\varphi(\alpha_X(\bar x'))\in\overset{\deleted{\mathord{-}}}{Y}(K_{k'})
\xleftarrow{\sim}\overset{\deleted{\mathord{-}}}{Y}(V_{k'})
\xleftarrow{\substack{\alpha_Y\\\sim}}\bar Y(k'),
\]
% slanted
\margin{Si $Y$ n'est pas propre sur $V$ et $\varphi$ \unsure{seulement} une application rationnelle, on a encore une ``application non partout définie'' $\bar X\longrightarrow\bar Y$, ayant les propriétés}
dont la localité est notre $\bar\varphi(x)$. On a donc par construction un hom
\[
k(\varphi(\bar x))\longrightarrow\widetilde{k(\bar x)}
\]
(où le $\sim$ dénote la clôture parfaite).

Pour la suite, $X,Y$ de type fini sur $K$, et \inserted{\unsure{séparés|sépr.|sérém.|sépar.}} intègres.

\uline{Proposition} 1\quad Supposons $\bar x\in\bar X$ \uline{épais} et $\varphi$ \uline{dominant \unsure{gén.\textsuperscript{t}|génér.\textsuperscript{t}} séparable} [i.e. fibre gén. séparable]. Alors\textsuperscript{\unsure{($\times$)|si}} $\bar y=\varphi(\bar x)$ \inserted{existe, et} est \uline{épais}, $k(\bar y)\subset{}k(\bar x)$, et l'application ponctuelle $\bar x_{\mathrm{réd}}\longrightarrow\bar y_{\mathrm{réd}}$ définie par cette inclusion provient d'un promorphisme d'indice fini sur un ouvert constructible dense $U\ni\bar x$, [d'image un \unsure{ens|un|ouvert} constructible de $\bar y_{\mathrm{réd}}$].
\scanpage{132}
% archive sheet number: 131
\uline{Corollaire}\quad Fonctorialité \tikz[baseline=(a.base)]{\node[draw,ellipse,inner sep=2pt] (a) {en $X_K$};} de l'ensemble $\operatorname{\text{Ép}}\deleted{\unsure{\mathrm{ais}}}(\bar X)$ des points \deleted{\unsure{de}} épais de $\bar X$ [munis de leurs corps résiduels, ----], pour des \deleted{\unsure{morphismes}}\inserted{applic. rationnelles} $X_K\unsure{\mathord{-}|\longrightarrow}Y_K$ \uline{dominantes gén.\textsuperscript{t} séparables de} $K$-\uline{schémas} \uline{propres}

\uline{Proposition} 2\quad Soit $S$ une partie constructible de $\bar X$, $X_K$-simple, [tel que $\unsure{l(N,x)|l(V,x)|(V,x)}$ y soit borné], et tel que $\bar\varphi(S)$ existe [p.ex. $Y$ propre sur $V$] et soit $Y_K$-simple. Alors $\bar\varphi(S)$ est constructible. De plus, il y a une partition \inserted{finie} $S=\coprod_i S_i$, les $S_i$ des sous-schémas irréductibles constructibles de $X$, \unsure{telles|telle} que $\varphi$ soit morphique d'indice fini sur chaque $S_i$, [et $\varphi(S_i)$ est alors un ens constructible irréductible \ldots]

\uline{Corollaire}\quad Fonctorialité de $X_K$ variant comme dessus, et de plus simple sur $K$, de l'ens $\operatorname{Cons}(\bar X)$ des parties constructibles de $\bar X$.
\scanpage{134}
% archive sheet number: 133
Considérons le foncteur d'anneaux [ou de préschémas \uline{quasi-compacts}]
\[
\Gamma:\quad A\leadsto A((t))^{*}=(A[[t]]_{t})^{*}
\]
\margin{($t$ une indéterminée)}

Ce foncteur est ind-représentable par un
\[
\mathscr G=\underset{n}{\varinjlim}\,(\mathscr G_n)
\]
\margin{\uline{$\varinjlim$} par abus de langage \unsure{ici !|où ?}}
$\mathscr G_n$ représente le ss foncteur $\Gamma_n$ de $\Gamma$, où une série formelle \deleted{$A\leadsto A((t))^{*}$} $\sum_{-\infty}^{+\infty}a_i t^i$ \inserted{ayant l'inverse $\sum_{-\infty}^{+\infty} b_j t^j$} est dans $\Gamma_n$ ssi $a_i=0,b_i=0$ pour $i<n$. $\mathscr G_n'$ est affine et en fait un sous-schéma fermé de $\operatorname{Spec}\mathbb Z[(a_i,b_i)_{i\geq-n}]$, [\unsure{fabriqué|défini} par les \deleted{relations} id\unsure{éal} des relations exprimant $(\sum_i a_i t^i)(\sum_j b_j t^j)=1$, savoir $a_{-n}b_{-n}=0,\ldots,$]

\deleted{Soit}\inserted{En particulier,} $\Gamma_0$ est le foncteur
\[
\Gamma_0:\quad A\leadsto A[[t]]^{*}
\]
qui est représenté par
\[
\mathscr G_0=\mathbb G_m\times\operatorname{Spec}\mathbb Z[(a_i)_{i\geq1}]
\]

Considérons le schéma
\[
\mathscr G'=\mathbb Z_{\mathbb Z}\times\mathscr G_0
\]
qui représente un foncteur ind-représentable
\[
\begin{aligned}
\unsure{\mathscr G'|\Gamma'}&=\underset{n}{\unsure{\underline{\lim}|\varinjlim}}\ \mathscr G'_n\\
\mathscr G'_n&=[-n,n]_{\mathbb Z}\times\mathscr G_0
\end{aligned}
\]
On a un morphisme canonique
\[
\mathscr G'\longrightarrow\mathscr G
\]
\scanpage{135}
% archive sheet number: 134
correspondant à des morphismes
\[
\mathscr G'_n\longrightarrow\mathscr G_n
\]
\deleted{\unsure{déf. à travers de foncteurs}} ou encore, composante par composante, à
\[
\varphi_n^{\unsure{\nu|v}}:\mathscr G_0\longrightarrow\mathscr G_n
\qquad \nu\in[-n,+n]
\]
donné \deleted{par} sur les foncteurs par
\[
\varphi_n^\nu\bigl(\deleted{\unsure{a|\Sigma}}{\textstyle\sum a_i t^i}_{i\geq0}\bigr)
=t^\nu\textstyle\sum a_i t^i
\]

Ceci posé, on \unsure{aura|a un|a}

\uline{Th}\quad (i) Le morph. $\varphi_n:\mathscr G'_n\unsure{\mathord{-}|\longrightarrow}\mathscr G_n$ est une \uline{immersion} \uline{fermée} \uline{surjective}
% slanted
\margin{donc $\mathscr G_n$ est réunion disjointe d'ouverts $\mathscr G_n^i$ $(i\in[-n,+n])$, d'ailleurs irréductibles, qui sont les images des $\varphi_n^i$}
(ii) Elle induit une immersion ouverte sur $\mathscr G_0^{\unsure{\times|\,}} \times\{\nu\}$ pour $\nu=\pm n$, mais pour $-n<\nu<+n$, le morphisme induit \uline{n'est pas} une imm ouverte [$\bar{\mathrm m}$ après chang\textsuperscript{t} de base quelconque $S\longrightarrow\operatorname{Spec}\mathbb Z$, $S\ne\varnothing$].
(iii) Soit $\mathscr G_n^i$ le sous schéma ouvert \inserted{induit} et fermé de $\mathscr G_n$ \deleted{\unsure{image}} sur l'image ensembliste \deleted{$\unsure{\varphi_n(\mathscr G}$}$\varphi_n(\mathscr G_0\times\{i\})$, où $|i|\leq n$.

Pour $n$ variable, on trouve ainsi un système inductif \inserted{d'immersions} de préschémas, ayant $\unsure{=|\simeq|\text{les}}$ \deleted{\unsure{ss-espace ponctuel|propriétés}}\inserted{\unsure{espace|mêmes}} top. sous-jacent. Pour \unsure{?} $n'\geq n\geq|i|$, le morphisme $\mathscr G_n^i\longrightarrow\mathscr G_{n'}^i$ n'est pas un isomorphisme, et ceci même sur une base $S\ne\varnothing$ quelc. [il suffit de regarder sur les corps premiers\ldots].
\scanpage{136}
% archive sheet number: 135
\diagram{\resizebox{\linewidth}{!}{\begin{tikzcd}[ampersand replacement=\&,column sep=small,row sep=small]
\& \operatorname{Pro}(\mathrm{Sch}_{/S}) \& \text{pro-préschémas}\\
\& \bigcup\ \substack{\text{sous-catégorie}\\\text{pleine}} \&\\
(\mathrm{Sch}_{/S}) \& \operatorname{Pro}'(\mathrm{Sch}_{/S})\arrow[l,"\varprojlim"'] \& \text{pro-préschémas} \& (X_i)_{i\in I},\quad\text{t.q.}\quad\exists i,\quad u_{ji}\ \text{affine pour }j\geq i\\
\& \bigcup\ \substack{\text{sous-catégorie}\\\text{pleine}} \&\\
\& \operatorname{Pro}''(\mathrm{Sch}_{/S})\arrow[uul,"\substack{\text{foncteur}\\\text{\unsure{plein\textsuperscript{mt}|plein|pleinement}}\\\text{fidèle}}"'] \& \text{pro-préschémas} \& {\begin{aligned}
&(X_i)_{i\in I}\quad\text{t.q.}\quad
\deleted{\left\{\begin{aligned}
&\unsure{X_i\ \text{présentation finie}}\\[-2pt]
&\exists i,\ \bigl[\unsure{X_i\ \text{quasi cpt}|X_i\ \text{qu.-compact et}}\ \text{de prés. finie}\\
&\quad u_{ji}\ \text{affine de prés. finie}
\end{aligned}\right.}\\
&\left\{\begin{aligned}
&\forall i\ \unsure{\mathrm{grd}|\mathrm{grand}},\ X_i\ \text{est qu. cpt et de présentation finie sur }S\\
&\exists i\ \text{t.q.}\ j\geq i\ \text{implique }u_{ji}\ \text{affine.}
\end{aligned}\right.
\end{aligned}}
\end{tikzcd}}}

Soient $\mathscr X\in\operatorname{Ob}\operatorname{Pro}''(\mathrm{Sch}_{/S})$, $\mathscr Y\in\operatorname{Ob}\operatorname{Pro}'(\mathrm{Sch}_{/S})$ \inserted{$\mathscr Y=(Y_j)$, les $Y_j$ \unsure{qu.cpts|quasi cpts},} alors
\[
\operatorname{Hom}_{\operatorname{Pro}(\mathrm{Sch}_{/S})}(\mathscr Y,\mathscr X)
\xrightarrow{\sim}\operatorname{Hom}(\varprojlim\mathscr Y,\varprojlim\mathscr X)
\]

\deleted{[En particulier, la catégorie}\inserted{Si \unsure{on|tout} est sur une base $S$}
\[
\operatorname{Hom}_{\operatorname{Pro}(\mathrm{Sch}_{/S})}(\mathscr Y,\mathscr X)
\simeq\operatorname{Hom}_S(\varprojlim\mathscr Y,\varprojlim\mathscr X)
\]

\deleted{Dans les provariétés de Serre}

\uline{Remarque}\quad Dans Serre et Néron, on utilise des systèmes projectifs de schémas \inserted{(algébriques)} qui sont \inserted{(pratiquement tjrs)} $\in\operatorname{Pro}''(\mathrm{Sch})$. \deleted{\unsure{Du|D}} Il y a avantage à les ``interpréter'' comme étant des véritables schémas sur $k$, \deleted{\unsure{de}} au lieu de pro-objets \ldots\ldots
\scanpage{138}
% archive sheet number: 137
\unsure{/|1}\quad \deleted{$\unsure{a_{i0}+|a_{i0}}$}\hfill $a_i,b_{ij},c_{ip}\in A$
\[
\left\{\begin{aligned}
f_i&=a_i\deleted{\unsure{a_{i0}}}+\sum_{1\leq j\leq n} b_{ij}t_j+\sum_{|p|\geq2}c_{ip}t^p\\
&1\leq i\leq n
\end{aligned}\right.
\qquad
\left\{\begin{aligned}
&a_i\in\mathfrak m\\
&\det(b_{ij})\notin\mathfrak m
\end{aligned}\right.
\]
\diagram{\begin{tikzcd}[ampersand replacement=\&,column sep=large]
(\lambda_1,\ldots\lambda_n)\in\mathfrak m^{(\times n)}
\arrow[r,maps to,decorate,decoration={snake,amplitude=.35pt,segment length=5pt},"\lambda_{(f=(f_1,\ldots f_n))}"']
\& (f_1(\lambda_1,\ldots,\lambda_n),\ldots,f_n(\lambda_1,\ldots\lambda_n))
\end{tikzcd}}

Si $A$ est d'ordre infinitésimal $N$, alors $\lambda_f$ est connu quand on connait $f\bmod (t_1,\ldots t_n)^{N+1}$, donc quand on connait \deleted{\unsure{les}} $\unsure{N|n}$ \uline{polynômes} \unsure{de|du|des} de degré $\leq N$.
\[
\unsure{\Longleftrightarrow|\text{ou}}\quad
\boxed{f_i(\lambda_1,\ldots\lambda_n)=a_i+\sum_j b_{ij}\lambda_j+
\sum_{2\leq|p|\leq N}c_{ip}\lambda_1^{p_1}\ldots\lambda_n^{p_n}}
\]
\diagram{\begin{tikzcd}[ampersand replacement=\&,row sep=1pt,column sep=small]
B=A[[t_1,\ldots t_n]] \& \& \mathfrak m^{\deleted{\unsure{1}}}\\
\& V_1\big[ \&\\
\& \& \mathfrak m^2\\
\& V_2\big[ \&\\
A\arrow[uuuu] \& \& \mathfrak m^3\\
\& \big[ \&\\
\& \big[ \& \unsure{\vdots|\,}\\
\& \& \mathfrak m^{N-1}\\
\& \big[ \&\\
\& \& \mathfrak m^N\unsure{\mathord{-}}\\
\& V_N\big[ \&\\
\& \& \mathfrak m^{N+1}=0
\end{tikzcd}}
\scanpage{139}
\mspage{- 132 -}
% archive sheet number: 138
\hfill n\textsuperscript{o} 316

terme constant, et utiliser le fait que la sous-algèbre des $X\handwritten{\in}\mathit{gl}(n,K)$ telles que $ {}^{t}XR+RX=0$ est de dimension $n(n-1)/2$ ).

c) Définir de même le \uline{groupe symplectique formel} à $2n$ variables sur un corps (commutatif) $K$ quelconque, et montrer que son algèbre de Lie s'identifie à la sous-algèbre de $\mathit{gl}(2n,K)$ formée des matrices $X$ telles que $ {}^{t}XA+AX=0$, où $A=\handwritten{\Bigl(}\begin{matrix}0&I_n\\-I_n&0\end{matrix}\handwritten{\Bigr)}$

27) - Soient $K$ un corps de caractéristique $p>0$, $G$ le groupe formel de dimension 2 défini par $f_1(x,y)=X_1+Y_1+X_1Y_1$, $f_2(x,y)=X_2+Y_2(1+X_1^p)$. Montrer que $G$ n'est pas commutatif, mais que son algèbre de Lie est commutative.
\scanpage{140}
% archive sheet number: 139
\[
\mathbb G_{n,N}=\underline{\operatorname{Aut}}_{\substack{\text{Algèbres}\\\text{augmentées}}}
\mathbb Z[t_1,\ldots,t_n]/(t_1,\ldots,t_n)^{N+1}
\]
$\Gamma_{n\deleted{\unsure{,N}}}=\operatorname{Grb}(\mathbb G_{n,N})$ groupe algébrique sur $k$ des changem\textsuperscript{ts} de coordonnées
\[
\left\{\begin{aligned}
&f_i(t_1,\ldots t_n)=\sum_j b_{ij}\unsure{t_j|t'_j}+
\sum_{2\leq|p|\leq N}c_{ip}t_1^{p_1}\ldots t_n^{p_n}\\
&b_{ij}\,c_{ip}\in A\qquad\det(b_{ij})\notin\mathfrak m
\end{aligned}\right.
\]

Ce groupe. opère sur $\mathfrak m^n$ si $A$ est d'ordre infinitésimal $\leq N$. \deleted{\unsure{Du|De|Donc}}
\[
\begin{aligned}
&g\in\Gamma_n \qquad x\leadsto T_f(x)\dot{+}a\\
&a\in\mathfrak m^n\\
&\qquad\deleted{\unsure{T_f|T_g}}\quad b'\dot{+}T_g(\deleted{\unsure{x|T}}a\dot{+}T_f(x))\\
&\qquad\qquad b\dot{+}T_g(a)\dot{+}
\end{aligned}
\]
\[
\begin{aligned}
&T_{\unsure{g_i|g'}}(a+b) &&g(x+y)=g(x)\\
&\deleted{\unsure{T_i}}\quad\deleted{\unsure{g_i|g}}\quad
\deleted{\unsure{f}}{\unsure{g|\deleted{g}}}(a+b)=g(a)+g'\\
&g(x+y)=
\end{aligned}
\]
\scanpage{141}
\mspage{- 129 -}
% archive sheet number: 140
\hfill n\textsuperscript{o} 316

où, dans la série $o_3$, tous les monômes sont de degré $\handwritten{\geqslant}3$, par rapport à l'ensemble des $Y_i$ et des $Z_i$. En déduire que :
\[
\begin{aligned}
u(f(x,y))-u(f(y,x))&=\handwritten{\sum}_{i\handwritten{<}j}(X_iY_j-X_jY_i)([T_i,T_j]^{(e)}(u))+\\
&\qquad +o'_3(x,y)
\end{aligned}\tag*{(3)}
\]
où tous les termes de $o'_3$ sont de degré total $\handwritten{\geqslant}3$. Montrer que si $G$ est commutatif, $g$ est commutative.

d) Soit $G'$ un second groupe formel, de dimension $n$, dont la loi de groupe est $f'$. On appelle \uline{homomorphisme formel} de $G$ dans $G'$ un système $F=(F_j(x))_{1\handwritten{\leqslant}j\handwritten{\leqslant}m}$ d'éléments de $A_x$ tel que l'on ait $f'(F(x),F(y))=F(f(x,y))$. Pour toute dérivation invariante à gauche $D\handwritten{\in}g$, montrer que $v\handwritten{\longrightarrow}D(v\unsure{o|\circ}F)$ est une dérivation invariante à gauche appartenant à l'algèbre de Lie $g'$ de $G'$. Si on la désigne par $F^{\handwritten{*}}(D)$, $F^*$ est un homomorphisme (un $p$-homomorphisme pour $p>0$) de $g$ dans $g'$. Si $(T_i)_{1\handwritten{\leqslant}i\handwritten{\leqslant}n}$, $(T'_j)_{1\handwritten{\leqslant}j\handwritten{\leqslant}m}$ sont les bases de $g$ et $g'$ telles que $T_i^{(e)}=D_i^{(e)}$ et $T'_j{}^{(e)}=D_j^{(e)}$, montrer que la matrice de $F^{\handwritten{*}}$, rapportée à ces bases, est la matrice $(\handwritten{\partial}F_j/\handwritten{\partial}X_i)_c$ (terme constant de $\handwritten{\partial}F_j/\handwritten{\partial}X_i$, où $i$ est l'indice des lignes, $j$ celui des colonnes).

25) - On appelle \uline{groupe linéaire formel} à $n$ variables sur le corps $K$, et on désigne par $\mathrm{GL}(n,K)$ (lorsqu'aucune confusion n'est possible) le groupe formel de dimension $n^2$ dont la loi de groupe est définie par $f_{ij}(u,v)=-\handwritten{\delta}_{ij}+\handwritten{\sum}_{k=1}^{n}(\handwritten{\delta}_{ik}+U_{ik})(\handwritten{\delta}_{kj}+V_{kj})$ ($\handwritten{\delta}_{ij}$ indice de Kronecker ; les $3n^2$ indéterminées de la définition générale de l'exerc. 24 sont ici notées $U_{ij}$, $V_{ij}$, $W_{ij}$ avec 2 indices variant de 1 à $n$). Montrer que, si $D_{ij}=\handwritten{\partial}/\handwritten{\partial}U_{ij}$, les dérivations
\scanpage{142}
% archive sheet number: 141
\deleted{$\unsure{\mathfrak m/\mathfrak m^2}$}
\[
\begin{aligned}
\text{espace affine sous}\qquad&\unsure{\mathscr T|\mathscr J|\mathscr V}^{0}_{B_0/A_0}\otimes V_1\\
\text{espace affine sous}\qquad&\mathscr T^{0}_{B_0/A_0}\otimes V_2\\
\text{\rule{4em}{0.3pt}}\qquad&\mathscr T^{0}_{B_0/A_0}\otimes V_3\qquad\mathscr T_{V_0/}
\end{aligned}
\]

$\mathbb G_{A,n}={}$ groupe algébrique sur $k$ des changem\unsure{\textsuperscript{ts}|ents} de base de \unsure{la|2} forme
\[
\left\{\begin{aligned}
&f_i=a_i+\sum_j b_{ij}t_j+\sum_{2\leq |p|\leq N}c_{ip}t^p
\qquad a_i,b_{ij},c_{ip}\in A\\
&a_i\in\mathfrak m,\quad\det(b_{ij})\notin\mathfrak m
\end{aligned}\right.
\]

$\mathbb G_n$ groupe sur $\mathbb Z$ des changem\textsuperscript{ts} de base
\[
\left\{\begin{aligned}
&f_i=a_i+\sum b_{ij}t_j+\sum_{2\leq |p|\leq N}c_{ip}t_1^{p_1}\ldots t_n^{p_n}\\
&\det(b_{ij})\ne0
\end{aligned}\right.
\]
i.e. des automorphismes de
\[
\deleted{\unsure{a_0+|a+t}}\qquad a_0\unsure{+|-}a_1t+a_2t^2
\]
\[
\begin{aligned}
&\deleted{a'_0+a'_1(a_0+a_1t)\unsure{+a'_2}}\ \unsure{\mathord{-}|=}\quad
a'_0+a'_1(a_0+a_1t+a_2t^2)+a'_2\bigl(\tikz[baseline=(q.base)]{\node[inner sep=0pt] (q) {$a_0+a_1t+a_2t^2$}; \draw ($(q.north west)!.14!(q.north east)$) -- ++(0,2pt) -- ($(q.north west)!.80!(q.north east)+(0,2pt)$) -- ++(0,-2pt); \draw ($(q.south west)!.17!(q.south east)$) -- ++(0,-2pt) -- ($(q.south west)!.37!(q.south east)+(0,-2pt)$) -- ++(0,2pt); \draw ($(q.south west)!.52!(q.south east)$) -- ++(0,-2pt) -- ($(q.south west)!.83!(q.south east)+(0,-2pt)$) -- ++(0,2pt); \draw ($(q.south west)!.14!(q.south east)$) -- ++(-1pt,-4pt); \draw ($(q.south west)!.42!(q.south east)$) -- ++(-1pt,-4pt);}\bigr)^2\\
&k[t]/(t^3)\qquad
\boxed{\begin{aligned}
&=(a'_0+a'_1a_0+a'_2a_0^2)\\
&\qquad +(a'_1a_1+\unsure{[|2}a'_2a_0a_1)t\\
&\qquad\qquad +(a'_1a_2+a'_2a_1^2+2a'_{\unsure{2|1}}a_0a_2)t^2
\end{aligned}}\\
&\hspace{8em}+2a'_2a_1a_2t^3\\
&\hspace{10em}+a'_2a_2^2t^4
\end{aligned}
\]
\scanpage{143}
\mspage{- 130 -}
% archive sheet number: 142
\hfill n\textsuperscript{o} 316

invariantes à gauche $X_{ij}$ telles que $X_{ij}^{(e)}=D_{ij}^{(e)}$ sont données par
\[
X_{ij}=(1+U_{ii})D_{ij}+\handwritten{\sum}_{k\ne i}U_{ki}D_{kj}.
\]
L'algèbre de Lie de $\mathrm{GL}(n,K)$ s'identifie à $\mathit{gl}(n,K)$ en identifiant $X_{ij}$ à élément $E_{ij}$ de la base canonique (utiliser la formule (3) de l'exercice.24).

Si $K$ est de caractéristique $p>0$, on a $X_{ii}^p=X_{ii}$ et $X_{ij}^p=0$ pour $i\ne j$.

\margin{\handwritten{$\unsure{\Pi|\pi}$}}
26) - a) Soient $K$ un corps de caractéristique $\ne2$, $E$ un espace vectoriel de dimension $n$ sur $K$, $\Phi$ une forme bilinéaire symétrique non dégénérée sur $E$. On suppose $E$ muni d'une base orthogonale pour $\Phi$ et on désigne par $R$ la matrice (diagonale) de $\Phi$ par rapport à cette base. On sait (\uline{Alg}. chap.IX, \S\quad, exe\unsure{r|f}c.\quad) que toute matrice (relative à la base considérée) $U$ du groupe orthogonal $G(\Phi)$ telle que $\det(I+U)\ne0$ s'écrit $U=(I-R^{-1}S)^{-1}(I+R^{-1}S)$, où $S=R(U-I)(U+I)^{-1}$ est une matrice alternée telle que $\det(I-R^{-1}S)\ne0$ ; inversement pour toute matrice $S$ satisfaisant à ces conditions, $U=(I-R^{-1}S)^{-1}(I+R^{-1}S)$ est une matrice de $G(\Phi)$ telle que $\det(I+U)\ne0$. Soient $S_{ij}$, $S'_{ij}$, $S''_{ij}$ $3n(n-1)/2$ indéterminées ($1\handwritten{\leqslant}i\handwritten{<}j\handwritten{\leqslant}n$) et soit $L$ un corps contenant l'anneau de séries formelles $K\handwritten{[}[S_{ij},S'_{ij},S''_{ij}]\handwritten{]}$. Si on désigne par $S$, $S'$, $S''$ les matrices $\handwritten{\sum}_{i\handwritten{<}j}S_{ij}(E_{ij}-E_{ji})$, $\handwritten{\sum}_{i\handwritten{<}j}S'_{ij}(E_{ij}-E_{ji})$, $\handwritten{\sum}_{i\handwritten{<}j}S''_{ij}(E_{ij}-E_{ji})$, on a $\det(I-R^{-1}S)\ne0$, et les analogues pour $S'$ et $S''$ ; si on pose $U=(I-R^{-1}S)^{-1}(I+R^{-1}S)$, $U'=(I-R^{-1}S')^{-1}(I+R^{-1}S')$, on a aussi $\det(I+UU')\ne0$. Posons $F(S,S')=R(UU'-I)(UU'+I)^{-1}=(f_{ij}(S,S'))$ ; les $f_{ij}(S,S')$ appartiennent à l'anneau des séries formelles
